%
\chapter[Procesos Qu\'{\i}micos en la Atm\'osfera]{Procesos qu\'{\i}micos en la atm\'osfera}

Los procesos qu\'{\i}micos en la atm\'osfera son el motor de transformaci\'on de los compuestos emitidos por fuentes naturales y actividad \gls{antropica}, determinando su impacto ambiental y ciclo de vida. Este cap\'{\i}tulo explora las reacciones clave que gobiernan la din\'amica de contaminantes en diferentes capas atmosf\'ericas, con \'enfasis en mecanismos de oxidaci\'on, interacciones foto-qu\'{\i}micas y transformaci\'on de especies cr\'{\i}ticas.

En la Secci\'on~\ref{7quim}, se analiza la \textbf{oxidaci\'on atmosf\'erica}, proceso fundamental que regula la degradaci\'on de contaminantes. Se describe la cadena de oxidaci\'on (\ref{71cox}), donde radicales libres como \ce{OH^.} act\'uan como ``detergentes atmosf\'ericos'', y la vida qu\'{\i}mica de los compuestos org\'anicos vol\'atiles (COVs) (\ref{72cov}), clave para entender su persistencia y formaci\'on de ozono troposf\'erico.

Las \textbf{Secciones \ref{qs02} y \ref{qno2}} profundizan en la qu\'{\i}mica de dos contaminantes cr\'{\i}ticos:
\begin{itemize}
    \item \textbf{Di\'oxido de azufre (\ce{SO2})}: Su conversi\'on a \'acido sulf\'urico (\ce{H2SO4}) mediante reacciones heterog\'eneas (\ref{qs02} ) y el papel de los compuestos reducidos de azufre (\ref{crso2}) en la nucleaci\'on de aerosoles.
    \item \textbf{Di\'oxido de nitr\'ogeno (\ce{NO2})}: Su rol en el ciclo del ozono (\ref{cozono}), actuando como precursor tanto en la formaci\'on de esmog como en la destrucci\'on estratosf\'erica de ozono (\ce{O3}).
\end{itemize}

Las Secciones~\ref{raminas} a  \ref{rchal} abordan la reactividad de grupos funcionales espec\'{\i}ficos:
\begin{itemize}
    \item \textbf{Aminas} (\ref{raminas}): Neutralizaci\'on de \'acidos en nubes y formaci\'on de nanopart\'{\i}culas
    \item \textbf{Compuestos de carbono} (\ref{recc}): Oxidaci\'on de metano (\ce{CH4}) y hidrocarburos arom\'aticos
    \item \textbf{Hal\'ogenos} (\ref{rchal}): Destrucci\'on catal\'{\i}tica de ozono estratosf\'erico (Ej.: \ce{Cl^.} de CFCs)
\end{itemize}

Estos procesos explican fen\'omenos globales como la lluvia \'acida, el agotamiento de la capa de ozono y el cambio clim\'atico, demostrando c\'omo interacciones moleculares escalan a impactos planetarios. El estudio integrado de estas reacciones provee herramientas para modelar escenarios de contaminaci\'on y dise\~nar estrategias de mitigaci\'on.

\section{Oxidaci\'on en la atm\'osfera}
\index{atmosfera@atm\'osfera!oxidacion@oxidaci\'on}\label{7quim}

La atm\'osfera es un ambiente oxidante: un hidrocarburo emitido puede oxidarse varias veces para formar los productos finales \ce{CO2} y \ce{H2O}. Los tres oxidantes principales son el radical hidroxilo\index{hidroxilo} (\ce{OH}), el ozono\index{ozono} (\ce{O3 }) y el nitrato\index{nitrato} (\ce{NO3}) (clasificados seg\'un su fuerza de oxidaci\'on decreciente).  Las reacciones de oxidaci\'on desempe\~nan un papel clave en la limpieza de la atm\'osfera (de lo contrario, las concentraciones de hidrocarburos ser\'{\i}an demasiado altas ).

En primera aproximaci\'on, la capacidad oxidante de la atm\'osfera viene dada por la concentraci\'on de OH (unas $10^6$ mol\'eculas cm${^{-3}}$ en la troposfera).

El radical hidroxilo OH a veces se denomina ``depurador qu\'{\i}mico', de la atm\'osfera. La mayor parte de los procesos de oxidaci\'on tienen lugar en las regiones tropicales.

El ozono (\ce{O3})  juega un papel decisivo para la producci\'on de OH a trav\'es de su fotolisis,

\reaction{ O3 + h\nu (\lambda \leqslant 310 nm) ->[L_1] O2 + O(^1D) \label{RA}}

El estado excitado de los \'atomos de ox\'{\i}geno, \ce{O(^1D)}, puede estabilizarse en la forma \ce{O(^3P)}, con la reacci\'on

\reaction{ O(^1D) + M ->[k_2] O(^3P) + M \label{RB}}

o producir OH con la reacci\'on

\reaction{ O(^1D) + H2O  ->[k_3]  2OH \label{RC}}
El punto clave es que s\'olo el estado fundamental excitado \ce{O(^1D)} reacciona con el vapor de agua (un compuesto estable). Como resultado, la producci\'on de OH est\'a directamente relacionada con la fotolisis del ozono.

La reacci\'on global resultante de estas tres reacciones es

\reaction{ O3 ->[h\nu, \ce{H2O}] 2OH \label{RD}}

Reglas que pueden ayudar y simplificar la forma de calcular los estados de oxidaci\'on\footnote{\textbf{Nota sobre oxidaci\'on}: Si un elemento pierde electrones, hay m\'as protones que electrones en el \'atomo, por lo que tendr\'a una carga positiva; por lo tanto, presentar\'a un estado de oxidaci\'on positivo. Si gana electrones, hay m\'as electrones que protones en el \'atomo, por lo que su carga y estado de oxidaci\'on ser\'an negativos.}:
\begin{enumerate}
\item El estado de oxidaci\'on de todos los elementos en su estado libre es 0. Esto se debe a que el elemento no ha perdido ni ganado electrones y, por tanto, es neutro. Este es el caso de los metales puros o las mol\'eculas biat\'omicas.

Por ejemplo: Zn, H y Cl.
\item La suma de los estados de oxidaci\'on de todos los \'atomos o iones de un compuesto neutro es igual a 0.

Por ejemplo: en el NaCl, el estado de oxidaci\'on del Na es +1 y el del Cl es -1. La suma de estos n\'umeros da como resultado 0.
\item La suma de los estados de oxidaci\'on de un ion es igual a su carga. Esto se aplica tanto a los iones monom\'ericos como a los complejos.

Por ejemplo: el estado de oxidaci\'on del ion monoat\'omico de Fl\'uor (\ce{F-}) es -1.

Por ejemplo: en el ion \ce{CO3^{-2}}, el C tiene un estado de oxidaci\'on de +4 y los tres O tienen un estado de oxidaci\'on de -2. Entonces, 4 + 3 (-2) = -2, que es la carga del ion.
\item En un ion o en un compuesto, el elemento m\'as electronegativo suele tener el estado de oxidaci\'on m\'as negativo. Tener presente que en la tabla peri\'odica la electronegatividad\index{electronegatividad} aumenta de izquierda a derecha en un per\'{\i}odo y disminuye de arriba hacia abajo en un grupo.

Por ejemplo: en el \ce{F2O}, el F es m\'as electronegativo que el ox\'{\i}geno, por lo que tiene el estado de oxidaci\'on m\'as negativo. Aqu\'{\i}, el F tiene un estado de oxidaci\'on de -1 y el O tiene un estado de oxidaci\'on de +2.
\end{enumerate}

\subsection{La cadena de oxidaci\'on}\index{oxidacion@oxidaci\'on!cadena de}\label{71cox}
Las reacciones m\'as frecuentes implican radicales con electrones libres en la capa de valencia externa (lo que implica que son altamente reactivos), como OH o \ce{HO2}.

Una cadena de oxidaci\'on se puede dividir en la siguiente secuencia de pasos sucesivos:

\begin{description}
\item[Iniciaci\'on ] Una especie estable (no es un radical; escrito como \ce{X_{norad}}) se disocia mediante una reacci\'on fotol\'{\i}tica, que conduce a la producci\'on de al menos un radical (\ce{X_{rad}}):
\reaction{X_{norad} +  h$\nu$  ->  X_{rad} \label{R1} }
La fotolisis del ozono (\ref{RA}) es un ejemplo de esta etapa.
\item[Propagaci\'on ] Los radicales resultantes reaccionan con especies estables (se mantiene la notaci\'on \ce{X_{norad}}, aunque no necesariamente sean las mismas especies que antes). Los nuevos radicales (tambi\'en escritos como \ce{X_{rad}}) se producen mediante reacciones de oxidaci\'on:
\reaction{X_{norad}  +  X_{rad}  -> X_{norad}  +  X_{rad}  \label{R2}}
\item[Reacci\'on de ramificaci\'on]  Las especies estables pueden disociarse durante una reacci\'on fotol\'{\i}tica y pueden generar otros radicales con reacciones similares a (\ref{R1}).
\item[Terminaci\'on ] Los radicales pueden reaccionar entre ellos para producir una especie estable, lo que finaliza la cadena de oxidaci\'on:

\reaction{X_{rad}  +  X_{rad}  -> X_{norad}    }
\end{description}

Las especies estables que se oxidan durante la etapa de propagaci\'on (\ref{R2}) deben considerarse como un ``combustible''  para la oxidaci\'on, ya que participan en la producci\'on de radicales. Normalmente es el caso de los COV, CO y \ce{CH4}.

\subsection{La vida qu\'{\i}mica de compuestos org\'anicos vol\'atiles}\label{72cov}

La vida qu\'{\i}mica de algunos compuestos org\'anicos vol\'atiles (COV)\index{COV} se muestra en el \textbf{Cuadro~\ref{tvcov}}. Se calculan con respecto a las especies oxidantes OH, \ce{O3} y \ce{NO3}, usando una ecuaci\'on similar a 
\begin{equation*}
\tau=\frac{1}{k_{\ce{OH},i}[\ce{OH}]}
\end{equation*}
(Como era de esperar, la vida qu\'{\i}mica aumenta cuando la concentraci\'on de OH disminuye y cuando $k_{\ce{OH},i}$ disminuye, es decir, cuando la especie es menos reactiva.)
\begin{table}[htp]
\caption[Vida Qu\'{\i}mica para algunos COV]{Vida Qu\'{\i}mica para algunos compuestos org\'anicos vol\'atiles  a \SI{298}{\kelvin} \footnotesize{(d para d\'{\i}a, h para hora y min para minuto)}}
\begin{center}
\begin{tabular}{|l|c|c|c|}\hline
Especie & Oxidaci\'on  &  Oxidaci\'on  &  Oxidaci\'on  \\
              &  con OH    & con \ce{O3} & con \ce{NO3} \\\hline
metano & 1837 d &  -- & -- \\
etano   & 48 d &  -- & 2960 d \\
butano & 4.8 d &  -- & 391 d \\
eteno & 1.4 d &  6.7 d & 107 d \\
propeno & 10.6 h & 1.1 d & 2.3 d\\
isopreno & 2.8 h & 20.2 h & 45 min \\
$\beta-$pineno & 3.5 h & 17.2 h & 12 min \\\hline
\end{tabular}
\end{center}
\label{tvcov}
\end{table}%
Las reacciones con OH suelen ser las m\'as r\'apidas y, como consecuencia, controlan la vida qu\'{\i}mica. Es f\'acil comprobar que el aumento de la ``complejidad qu\'{\i}mica'' da como resultado un aumento de la reactividad y, por tanto, una disminuci\'on de la vida \'util.

\section{Qu\'{\i}mica del di\'oxido de azufre}\label{qs02}
\index{SO2@\ce{SO2}!quimica@qu\'{\i}mica}
Desde le punto termodin\'amico, el di\'oxido de azufre (\ce{SO2}) posee una fuerte tendencia a reaccionar con el ox\'{\i}geno del aire.
\reaction{2SO2 +O2 -> 2SO3 \label{RS1}}

La proporci\'on de concentraciones en equilibrio de [\ce{SO3}]/[\ce{SO2}] \index{SO2@\ce{SO2}} \index{azufre!oxidos@\'oxidos} es cerca de $8\times10^{11}$ en aire a 1 atm\'osfera y \SI{25}{\celsius}. Sin embargo la rapidez de reacci\'on de la reacci\'on~\ref{RS1}  es muy baja en fase gas en una atm\'osfera libre de catalizadores por lo que se puede descartar como fuente de \ce{SO3}. Una vez formado el \ce{SO3}  reacciona muy r\'apidamente con el vapor de agua par formar el \'acido sulf\'urico 

\reaction{ SO3  + H2O -> H2SO4(ac) \label{RS2}}
cualquier proceso en el que el \ce{SO3} se forme en una atm\'osfera h\'umeda pude considerarse equivalente a la formaci\'on del \ce{H2SO4}.

Para explica la oxidaci\'on del \ce{SO2}, uno debe ver las reacciones diferentes a la directa. El di\'oxido de azufre absorbe la luz en la regi\'on ultravioleta de  la luz solar incidente en la troposfera, llevando a un estado excitado a la mol\'ecula de \ce{SO2}.
\reaction{ SO2  + h$\nu$ -> SO2^* }
El estado excitado del \ce{SO2^*}, est\'a no disociado. Solo longitudes de onda debajo de los 218 nm, los cuales no penetran a la troposfera, proveen la energ\'{\i}a suficiente para permitir la foto-disociaci\'on de \ce{SO2}. Si cada mol\'ecula que \ce{SO2} que queda foto-excitada se oxidase mediante la reacci\'on con el \ce{O2} u otra especie, se tendr\'{\i}a que la vida del \ce{SO2} es tan baja como de 52 min. Sin embargo, conocemos que el tiempo de vida del \ce{SO2} es mucho mayor que eso.

El  \ce{SO2} reacciona en condiciones de la troposfera v\'{\i}a procesos fase gas y acuosa y se remueve f\'{\i}sicamente v\'{\i}a deposito seco y h\'umedo. Con respecto a la reacci\'on fase gas, la reacci\'on con el OH es la dominante:
\reaction{ SO2  + OH^. +  M ->  HOSO2^. + M }
seguido de la regeneraci\'on del radical \ce{HO2}\index{radical!HO2@\ce{HO2}}
\reaction{  HOSO2^. + O2 -> HO2^.  + SO3 }
\reaction{  HO2^. + NO -> NO2  + OH^. }
y el tri\'oxido se convierte a \'acido sulf\'urico v\'{\i}a la reacci\'on \ref{RS2}:
\reaction*{ SO3  + H2O -> H2SO4} 
El tiempo de vida del \ce{SO2} basado en la reacci\'on con el radical OH, a concentraciones t\'{\i}picas atmosf\'ericas, es de cerca a una semana. El \ce{SO2} es uno de los gases que se remueve de forma razonablemente eficiente de la atm\'osfera por dep\'osito seco.

\subsection{Compuestos reducidos de azufre}\index{azufre!compuestos reducidos}\label{crso2}
Las especies reducidas de azufre  RSR' reacciona con los radicales OH y \ce{NO3}. \'Estas son emitidas por procesos biog\'enicos liberando especies como monosulfuro de carbono (CS), \'acido sulfh\'{\i}drico (\ce{H2S}), sulfuro de carbonilo (\gls{COS}), dimetil sulfuro (\ce{CH3SCH3} o \gls{DMS}) y bisulfuro de metilo (\ce{CH3SSCH3}). Para el \ce{H2S}, el proceso troposf\'erico dominante de remoci\'on involucra a la reacci\'on del radical OH,

\reaction{ H2S  + OH^.  ->  SH^. + H2O }

El tiempo de vida del \ce{H2S} por esta reacci\'on es cerca de 70 horas. Parece que el radical \ce{SH^.} sigue una serie de reacciones resultado en la formaci\'on de  \ce{SO2}. Para el metanotiol (\ce{CH3SH})\index{metanotiol} las reacciones tanto con el OH y \ce{NO3} conducen al grupo \ce{CH3S(OH)H}, que se cree que pasa, a trav\'es del \ce{CH3S}, a formaldeh\'{\i}do (HCHO),  \ce{SO2} y al \'acido metasulf\'onico (\ce{CH3SO3H})\index{acido@\'acido!metasulfonico@metasulf\'onico}.

El sulfuro de dimetilo (DMS), \ce{CH3SCH3}\index{DMS@\textbf{DMS}}, es un \gls{tioeter}, es el mayor contribuidor natural al flujo global de azufre. La vida del DMS en atm\'osferas marinas es el resultado de ambas reacciones con el OH y \ce{NO3} estando en el borde de uno a varios d\'{\i}as, con la mayor\'{\i}a de las rutas ocurriendo con el OH a bajas latitudes y con el \ce{NO3} en regiones fr\'{\i}as  y obscuras. Debido a que la fuente fotoqu\'{\i}mica del OH, la remoci\'on del DMS por OH se da durante el d\'{\i}a, provocando un ciclo reducido la concentraci\'on del DMS.


\section{Qu\'{\i}mica del di\'oxido de nitr\'ogeno}\label{qno2}
\index{NO2@\ce{NO2}!quimica@qu\'{\i}mica}
Las especies predominantes de nitr\'ogeno en la qu\'{\i}mica de atm\'osferas naturales y contaminadas son el NO, \ce{NO2} y \ce{HNO3}. 
Cuando el NO y \ce{NO2}  est\'an en presencia de la luz solar, se da la formaci\'on de ozono como resultado de la fotolisis del \ce{NO2}  en longitudes de onda ($\lambda$) <420 nm,

\reaction{ NO2 + h$\nu$  ->[k_1] NO + O^. \label{RN1}}

\reaction{O^. + O2 + M ->[k_2] O3 + M      \label{RN2}}

donde M representa \ce{O2} o \ce{N2} u otro tipo de tercer mol\'ecula que absorbe el exceso de energ\'{\i}a vibracional y de este modo estabiliza la mol\'ecula de \ce{O3}. No existen fuentes significativas de ozono en la atm\'osfera otras que la reacci\'on \ref{RN2}. Una vez formado, el \ce{O3} reacciona con el NO para regenerar el \ce{NO2}.

\reaction{O3 + NO  ->[k_3]  NO2 + O2       \label{RN3}}

Consideremos un momento la din\'amica del sistema en el cual s\'olo estas tres reacciones se dan. Asumiendo que se conocen las concentraciones iniciales de NO y \ce{NO2}, $[\ce{NO}]_0$ y $[\ce{NO2}]_0$,en el aire en un reactor a volumen y presi\'on constante e irradiado. La tasa de cambio de la concentraci\'on de \ce{NO2}, despu\'es de que la irradiaci\'on inicio est\'a dada por

\begin{equation}
\frac{\mathrm{d}[\ce{NO2}]}{\mathrm{d}t}=-k_1[\ce{NO2}]+k_3[\ce{O3}][\ce{NO}]
\label{RN4}
\end{equation}

considerando el $[\ce{O2}]$ como constante hay cuatro especies en el sistema: \ce{NO2}, NO, O y \ce{O3}. Podemos escribir la ecuaciones par la dinamic\'e del NO, O, y \ce{O3} como se realiz\'o para el \ce{NO2}. Por ejemplo, la ecuaci\'on para el  [O] es:

\begin{equation}
\frac{\mathrm{d}[\ce{O}]}{\mathrm{d}t}=k_1[\ce{NO2}]-k_2[\ce{O}][\ce{O2}][\ce{M}]
\label{RN5}
\end{equation}

Sin embargo, si se desea evaluar el lado derecho num\'ericamente veremos que es muy cercano a cero. F\'{\i}sicamente, esto significa que el \'atomo de ox\'{\i}geno es tan reactivo que desaparece por la reacci\'on \ref{RN2}  tan r\'apido como se ha formado en la reacci\'on \ref{RN1}. Para tratar con especies altamente reactivas como el \'atomo de ox\'{\i}geno es acostumbrado invocar la aproximaci\'on de estado seudo estable (PSSA) y supongamos por tanto que la tasa de formaci\'on es exactamente igual a la tasa de desaparici\'on, por ejemplo

\begin{equation*}
k_1[\ce{NO2}]=k_2[\ce{O}][\ce{O2}][\ce{M}]
\end{equation*}

En el estado estacionario la concentraci\'on del \'atomo de ox\'{\i}geno en este sistema esta dado por

\begin{equation}
[\ce{O}]=\frac{k_1[\ce{NO2}]}{k_2[\ce{O2}][\ce{M}]}
\label{RN7}
\end{equation}

Notar que $[\textrm{O}]$ no es constante; varia con el $[\ce{NO2}]$ en tal forma que en cualquier instante el balance se alcanza entre la tasa de producci\'on y p\'erdida. Lo que esta aproximaci\'on significa que que la concentraci\'on del \'atomo de ox\'{\i}geno se ajusta a cambios en la concentraci\'on de \ce{NO2} en varios ordenes de magnitud mas r\'apido que el cambio en la concentraci\'on de NO2. As\'{\i}, en la escala de tiempo de la din\'amica del \ce{NO2}, $[\textrm{O}]$ siempre satisface \ref{RN7}.

Sin embargo, de \ref{RN4} y \ref{RN5} vemos que de esas tres reacciones alcanzar\'an un punto donde el \ce{NO2} destruido y reformado tan r\'apido que se mantiene el ciclo del estado estacionario. El estado estacionario de la concentraci\'on de ozono esta dado por

\begin{equation}
[\ce{O3}]=\frac{k_1[\ce{NO2}]}{k_3[\ce{NO}]}
\label{RN8}
\end{equation}

esta expresi\'on, result\'o del an\'alisis del estado estacionario de las reacciones \ref{RN1} y \ref{RN3}, ha sido nombrada como la relaci\'on del estado fotoestacionario\index{estado fotoestacionario}. Observemos que la concentraci\'on de ozono en el estado estacionario es proporcional a la relaci\'on $[\ce{NO2}]/[\ce{NO}]$. Para calcular  $[\ce{NO2}]$ y $[\ce{NO}]$. Se obtienen de la conservaci\'on del nitr\'ogeno
\begin{equation*}
[\ce{NO}] + [\ce{NO2}] = [\ce{NO}]_0 + [\ce{NO2}]_0
\end{equation*}

y de la reacci\'on estequiom\'etrica del \ce{O3} con el \ce{NO}

\begin{equation*}
[\ce{O3}]_0 - [\ce{O3}] = [\ce{NO}]_0 - [\ce{NO}]
\end{equation*}

Resolviendo para $[\ce{O3}]$, se obtiene que la reacci\'on para la concentraci\'on de ozono formada en el estado estacionario por la irradiaci\'on de cualquier mezcla de NO, \ce{NO2}, \ce{O3} y exceso de \ce{O2}
\begin{eqnarray}
 [\ce{O3}]&=&-\frac{1}{2}\left([\ce{NO}]_0-[\ce{O3}]_0+\frac{k_1}{k_3}\right) \nonumber\\
  &+&\frac{1}{2}\left[\left([\ce{NO}]_0-[\ce{O3}]_0+\frac{k_1}{k_3}\right)^2+\frac{4k_1}{k_3} \left([\ce{NO2}]_0+[\ce{O3}]_0\right)\right]^{1/2}
\label{RN9}
\end{eqnarray}

Si $[\ce{O3}]_0=[\ce{NO}]_0=0$ la reacci\'on \ref{RN9} se reduce a 
\begin{equation}
 [\ce{O3}]=\frac{1}{2}\left(\left[\frac{k_1}{k_3}\right]^2+\frac{4k_1}{k_3} [\ce{NO2}]_0\right)^{1/2} -\frac{k_1}{2k_3}
\label{RN10}
\end{equation}

un valor t\'{\i}pico de $k_1/k_3$ expresado en unidades de raz\'on de mezcla es 10 ppb, as\'{\i} podemos calcular la concentraci\'on de ozono inicial de \ce{NO2} con   $[\ce{O3}]_0=[\ce{NO}]_0=0$.

\subsection{Concentraci\'on de ozono}\label{cozono}
Si $[\ce{NO2}]_0=0$, entonces $[\ce{O3}]=0$. Esto es claro ya que sin \ce{NO2} no se produce el ox\'igeno at\'omico y por lo tanto el ozono.  As\'{\i} la concentraci\'on m\'axima de ozono en el estado estacionario se alcanzara con una concentraci\'on inicial pura de \ce{NO2} \textbf{Cuadro~\ref{ozono_foto}}. Las concentraciones de ozono en atm\'osferas urbanas y regionales com\'unmente son mayores que las obtenidas en este c\'alculo (120 ppb). Como la mayor parte del  \ce{NO$_x$} es emitido en la forma de NO y no de  \ce{NO2}, la concentraci\'on de ozono alcanzada, si solo se emplean las reacciones \ref{RN1} y \ref{RN2} serian muy bajas comparadas con  las concentraciones observadas. Esto concluye que debe existir otras reacciones a \ref{RN1} y \ref{RN3} que son importantes en el aire troposf\'erico en donde concentraciones relativamente altas se dan.
\begin{table}[htp]
\caption[Concentraci\'on de ozono sin COV]{Concentraci\'on de ozono considerando reacciones de estado fotoestacionario.}
\begin{center}
\begin{tabular}{|c|c|}\hline
$[\ce{NO2}]_0$ (ppb) &$[\ce{O3}]_0$ (ppb) \\\hline
100   & 27\\
1000 & 95 \\ \hline
\end{tabular}
\end{center}
\label{ozono_foto}
\end{table}%

\section{Reacciones de las aminas}\index{aminas}\label{raminas}
Las principales reacciones atmosf\'ericas en fase gaseosa de las aminas involucran el radical OH:
\reaction*{NH3 + OH ->H2O + NH2^.}
\reaction*{CH3NH2 + OH -> H2O + CH3NH^.}
\reaction*{(CH3)2NH  + OH -> H2O + (CH3)2N^.}
La aminas tambi\'en reacciona con el \'acido n\'{\i}trico gaseoso para formar las sales de nitrato correspondientes.
La reacci\'on de amon\'{\i}aco con el vapor de \'acido n\'{\i}trico da el nitrato de amonio que es la ruta principal de formaci\'on de nitrato part\'{\i}cula
\reaction{NH3 + HNO3 ->  NH4NO3_(s) }
El amon\'{\i}aco, es muy soluble y reactivo con los \'acidos atmosf\'ericos, se remueve preferentemente por aquellas rutas que mediante la reacci\'on con el radical OH, la cual es relativamente lenta.

\section{Reacciones de los compuestos de carbono}\label{recc}
El mon\'oxido de carbono (CO)\index{monoxido@mon\'oxido!de carbono}, es el producto principal de los procesos de combusti\'on tanto de motores, generaci\'on de electricidad e incineraci\'on. Su rol como contaminantes es debido a su alta toxicidad y es importante en la producci\'on de \'atomos de hidr\'ogeno y la formaci\'on de radicales hidroxiperoxi\index{radical!hidroperoxi}. Sin embargo los compuesto se carbono van m\'as all\'a del mon\'oxido de carbono y tambi\'en contribuyen de manera importante a la contaminaci\'on del aire. Hidrocarburos sin quemar de las emisiones vehiculares, solventes industriales, afluentes de plantas industriales, etc., todos contribuyen a la interacci\'on con los productos de otras reacciones fotoinducidas.

Los tres principales mecanismos de reacci\'on de compuesto de carbono en la atm\'osfera involucran el ox\'{\i}geno singelete, radicales hidroxilo y ozono.
Todos ellos reacciona con las mol\'eculas org\'anicas para formar radicales org\'anicos. Esos radicales son precursores, entre otras cosas, de oxidantes conteniendo carb\'on que, como el ozono y \'acido n\'{\i}trico, son contaminantes que afectan de forma importante a las plantas y materiales. Tambi\'en est\'an los compuestos que contiene carb\'on del esmog que inducen la irritaci\'on de ojos y garganta.
\begin{description}
\item[Ox\'{\i}geno singelete ] \index{singelete} Este se produce a partir de la fotodisociaci\'on de una mol\'ecula de ox\'{\i}geno, reacciona con los hidrocarburos para producir radicales org\'anicos de la siguiente forma:
\reaction*{RH  +  O(^1D) ->  R^.  + HO^.}
\item[Radical Hidroxilo] \index{radical!hidroxilo} Los radicales hidroxilo, este tambi\'en producido de la reacci\'on del agua con el ox\'{\i}geno singelete, tambi\'en reacciona con hidrocarburos
\reaction*{RH  +   HO^.  ->  R^.  + H2O}
para producir radicales org\'anicos\index{radical!organico@org\'anico}.
Los radicales org\'anicos despu\'es reaccionan r\'apidamente con el ox\'{\i}geno para formar radicales peroxi:

\reaction*{ R^.  + O2 -> R-O-O^.  }
los cuales se convierten, reaccionando con el \ce{NO2} para formal alquilperoxinitratos\index{alquilperoxinitrato}
\reaction*{ R-O-O^.  + NO2  ->  R-O-ONO2 }
que son un tipo de contaminantes atmosf\'ericos oxidantes.
\item[Ozono] El mecanismo de reacci\'on del ozono en la formaci\'on de radicales es mas complicada. El ozono no extrae hidr\'ogenos de los compuestos org\'anicos. Mas com\'un, forma radicales mediante la creaci\'on con hidrocarburos instaurados como el etileno o propileno, formando compuestos adicionados con ox\'{\i}geno de acuerdo con
\begin{figure}[htbp]
\begin{center}
\begin{picture}(55,18)
% ozono
\put(2,4){\ce{O3} + \ce{H2C\bond{2}CH2 -> } }
\put(35,2){\ce{CH2\bond{-}CH2}}
%
\put(35,5){\line(-1,2){2}}
\put(46,5){\line(1,2){2}}
% oxigenos
\put(32,10){\ce{O}}
\put(47,10){\ce{O}}
%
\put(35,13){\line(5,3){4}}
\put(46,13){\line(-5,3){4}}
%
\put(39,15){\ce{O}}
\end{picture}
\caption{Reacci\'on del ozono (\ce{O3}) con etileno (\ce{CH2\bond{=}CH2}).}
\label{03-etileno}
\end{center}
\end{figure}

los cuales se descomponen para formar aldeh\'{\i}dos o diradicales inestables de la forma:
\begin{center}
\begin{picture}(50,14 )
% ozono
\put(2,5){\ce{H-C}  }
%
\put(10,8){\line(5,3){4}}
\put(10,5){\line(5,-3){4}}
\put(10.5,6){\line(5,-3){4}}
%
\put(15,10){\ce{H}  }
\put(15, 1){\ce{O}  }
\put(21,5){y}
\put(26,3){\ce{H3C-C-O-O^.}}
\put(36,0.2){\ce{^.}}
\put(36,6){\line(0,1){2.5}}
\put(35,9){\ce{H}}
\end{picture}
\end{center}

Los diradicales, se descomponen en otros radicales, como por ejemplo:
\reaction*{HO^.  \textrm{  y  } H3C-CO^.}
Esos y otros radicales reacciona con otras mol\'eculas org\'anicas en numerosas formas y en un sinf\'{\i}n de otros tipos de carbonilos. 

Los aldeh\'{\i}dos formados en las reacciones con ozono tambi\'en contribuyen la contaminaci\'on atmosf\'erica sirviendo como precursores de otra clase de compuestos org\'anicos oxigenados llamados nitratos de peroxiacil. Un ejemplo de uno de esto es el nitrato de peroxiacetil (PAN)\index{PAN}, tambi\'en conocido como lacrim\'ogeno o irritante de ojos, entornado en el esmog fotoqu\'{\i}mico.

El PAN se forma como consecuencia de la reacci\'on de aldeh\'{\i}dos con el radical hidroxilo:
\begin{center}
\begin{picture}(60,14)

% ozono
\put(2,5){\ce{H-C}  }
%
\put(10,8){\line(5,3){4}}
\put(10,5){\line(5,-3){4}}
\put(10.5,6){\line(5,-3){4}}
%
\put(15,10){\ce{H}  }
\put(15, 1){\ce{O}  }
% continua la reaccion
\put(19,5){\ce{+  HO^. -> H-C^. + H2O}}
\put(44.5,8){\line(0,1){2.5}}
\put(45.5,8){\line(0,1){2.5}}
\put(43.7,11){\ce{O}  }
\end{picture}
\end{center}
seguido de la reacci\'on con ox\'{\i}geno:
\begin{center}
\begin{picture}(55,17)
% reaccion con oxigeno
\put(2,8){\ce{H-C^. + O2 -> H-C-O-O^.}  }
\put(8,11){\line(0,1){2.5}}
\put(9,11){\line(0,1){2.5}}
\put(7.2,14){\ce{O}  }
%
\put(34,11){\line(0,1){2.5}}
\put(35,11){\line(0,1){2.5}}
\put(33.2,14){\ce{O}  }
\end{picture}
\end{center}
para formar el radical peroxiacilo. Y el radical peroxiacilo, en turno, reacciona con el \ce{NO2} para formar PAN
\begin{center}
\begin{picture}(70,11)
%\put(0,0){\line(1,0){70}}
%\put(0,11){\line(1,0){70}}
%\put(0,0){\line(0,1){11}}
% reaccion con oxigeno
\put(2,2){\ce{ H-C-O-O^. + NO2 -> H-C-O-O-NO2}  }
\put(8,5){\line(0,1){2.5}}
\put(9,5){\line(0,1){2.5}}
\put(7.2,8){\ce{O}  }
\put(47,5){\line(0,1){2.5}}
\put(48,5){\line(0,1){2.5}}
\put(46.2,8){\ce{O}  }
\end{picture}
\end{center}
Existe una multitud de otras reacciones inducidas por la fotoqu\'{\i}mica que se ha reconocido como contribuidoras a la contaminaci\'on del aire y a la formaci\'on del esmog.
\end{description}

\section{Reacciones de los compuestos con hal\'ogenos}
\index{halogeno@hal\'ogeno!reacciones}\label{rchal}

Una vez en la atm\'osfera las mol\'eculas org\'anicas halogenadas se descomponen por la fot\'olisis directa o por el ataque de los radiales hydroxyl. Ambas rutas conducen a la liberaci\'on del \'atomo de hal\'ogeno. Por ejemplo, la reacci\'on del OH con el cloruro de metilo conduce a:
\begin{align*}
\ce{CH3Cl + OH^. &->  ^.CH2Cl + H2O \\
      ^.CH2Cl + O2 & ->[M] CH2ClOO^. \\
      CH2ClOO^. + NO & -> NO2 + HCHO + Cl^. }
\end{align*}
Los \'atomos de los hal\'ogenos son muy reactivos con los hidrocarburos, llevando a la formaci\'on de haluros de hidr\'ogeno mediante la abstracci\'on del hidr\'ogeno 
\reaction{Cl^. + RH -> R + HCl \label{RH1} }
El fl\'uor y cloro reaccionar f\'acilmente de esta manera. Los \'atomos de Br pueden extraer un \'atomo de hidr\'ogeno solo del radical  \ce{HO2} o de los aldeh\'{\i}dos. Los \'atomos de yodo son menos reactivos. La alternativa de la reacci\'on anterior es la oxidaci\'on del \'atomo de hal\'ogeno (X=Cl, Br, I) por ozono:
\reaction{X + O3  -> XO^. + O2  \label{RH2} }
Debido a la reactividad decreciente de los hal\'ogenos, mientras que van de F a I, hacia compuestos que contiene H de la forma HX, la fracci\'on de \'atomos de hal\'ogenos libres que reaccionan en la ecuaci\'on \ref{RH2}, en contraposici\'on a \ref{RH1}, es F$\sim$0\%,Cl, 50\%, Br, 99\%, I, 100\%.

El haluro de hal\'ogeno, HX, puede reaccionar con el OH,
\reaction{HX + OH  -> X + H2O  \label{RH3} }
el cual regresa un \'atomo de hal\'ogeno X a la reserva de hal\'ogenos. Los radicales de \'oxido de hal\'ogeno puede seguir varias reacciones. Esas incluyen la fot\'olisis (importante para X=I, Br y en menor medida Cl).
\reaction*{XO + h$\nu$  -> X + O^.   }
reactivada con NO,
\reaction*{XO + NO  ->  X + NO2   }
y reacci\'on con \ce{HO2}
\reaction*{XO + HO2  ->  XOH + O2   }
En las masas de aire influenciadas por las emisiones antr\'opicas (la concentraci\'on de NO en el orden de 1 ppm), las razones [XO]/[X] son del orden de 10 a 100 (X=Cl) y a la unidad (X=Br, I).

Las reacciones con los \'oxidos de nitr\'ogeno \ce{NO2} o \ce{N2O5} con NaX contenidas en el aerosol salino del mar pueden conducir a la formaci\'on de XNO o \ce{XNO2}, respectivamente,
\reaction*{N2O5 + NaX_{(s)} ->  XNO2 + NaNO3_{(s)}}

Los harulos de hidr\'ogeno se pueden liberar del aerosol salino marino por medio de la reacci\'on con \'acidos fuertes como el \ce{H2SO4} y \ce{HNO3}:
\reaction*{H2SO4 + 2NaX_{(s)} -> 2HX + Na2SO4_{(s)} }
La raz\'on es que debido a que las constantes de velocidad de las reacciones de los \'atomos de hal\'ogeno, particularmente Cl, con hidrocarburos son significativamente mayores que las de las correspondientes \ce{OH + HC}, los hidrocarburos se pueden eliminar eficazmente mediante la reacci\'on con \'atomos de hal\'ogeno.

En la \textbf{Figura~\ref{halo_01}} se muestra el intercambio de los HCFC y sus derivados en la atm\'osfera algunos de los procesos de transporte puede durar algunos d\'{\i}as y otros llevan a\~nos.

\begin{figure}[htbp]
\begin{center}
\includegraphics[width=0.95\textwidth]{halogenos_01.png}
\caption{Intercambio de hal\'ogenos en la atm\'osfera.}
\label{halo_01}
\end{center}
\end{figure}

% =============================================================================
% exercises/ejercicios_cap08.tex
% Ejercicios propuestos — Capítulo 7: Procesos Químicos en la Atmósfera
% =============================================================================

\section*{Ejercicios Propuestos}
\addcontentsline{toc}{section}{Ejercicios propuestos}
\label{sec:ejercicios_cap08}

\begin{bloque_ejercicios}[Ejercicios --- Cap\'{\i}tulo 7: Procesos Qu\'{\i}micos en la Atm\'osfera]

\begin{ejercicios}

% ----- Ejercicio 1: Oxidantes atmosféricos y cadena de oxidación ------------
\item Los tres principales oxidantes en la atm\'osfera son el radical
hidroxilo (\ce{OH^.}), el ozono (\ce{O3}) y el radical nitrato (\ce{NO3}).

\begin{enumerate}[label=\alph*)]
    \item Ordene los tres oxidantes de \textbf{mayor a menor} fuerza
          oxidante seg\'un se indica en el cap\'{\i}tulo e identifique en
          qu\'e per\'{\i}odo del d\'{\i}a (diurno o nocturno) domina
          cada uno como agente de remoci\'on de COV. Justifique.

    \item La producci\'on de \ce{OH^.} en la troposfera depende de la
          fotolisis del ozono. Escriba la secuencia de tres reacciones
          (denominadas \ref{RA}, \ref{RB} y \ref{RC} en el cap\'{\i}tulo)
          que llevan desde \ce{O3} hasta \ce{2OH^.}, e identifique
          en cu\'al de ellas interviene el vapor de agua.

    \item Usando la definici\'on de vida qu\'{\i}mica
          $\tau = 1/\left(k_{\ce{OH},i}\,[\ce{OH}]\right)$,
          explique por qu\'e el metano (\ce{CH4}, $\tau = 1837$~d)
          tiene una vida qu\'{\i}mica much\'{\i}simo mayor que
          el isopreno ($\tau = 2.8$~h) frente al mismo oxidante.
          ?`Qu\'e implica esto para su transporte global?

    \item Clasifique las cuatro etapas de una cadena de oxidaci\'on
          (iniciaci\'on, propagaci\'on, ramificaci\'on y terminaci\'on)
          e indique con un ejemplo tomado del cap\'{\i}tulo a qu\'e etapa
          corresponde la fotolisis del ozono (\ref{RA}).
\end{enumerate}

% ----- Ejercicio 2: Vida química de COV ------------------------------------
\item Con base en el \textbf{Cuadro~\ref{tvcov}} de vidas qu\'{\i}micas
de COV a \SI{298}{\kelvin}:

\begin{enumerate}[label=\alph*)]
    \item Reproduzca y complete la siguiente tabla indicando el oxidante
          \textbf{dominante} (el que produce la vida m\'as corta) para
          cada especie:

{%
\centering
\begin{tabular}{lcccl}
\toprule
\textbf{COV} & $\tau_{\ce{OH}}$ & $\tau_{\ce{O3}}$ &
$\tau_{\ce{NO3}}$ & \textbf{Oxidante dominante} \\
\midrule
Metano    & 1837 d & ---    & ---    & \\
Eteno     & 1.4 d  & 6.7 d  & 107 d  & \\
Propeno   & 10.6 h & 1.1 d  & 2.3 d  & \\
Isopreno  & 2.8 h  & 20.2 h & 45 min & \\
$\beta$-pineno & 3.5 h & 17.2 h & 12 min & \\
\bottomrule
\end{tabular}
\par
}%

    \item Para el propeno, calcule la constante de velocidad de segundo
          orden $k_{\ce{OH},\text{propeno}}$ sabiendo que la
          concentraci\'on t\'{\i}pica de \ce{OH^.} en la troposfera
          es $[\ce{OH}] = 1\times10^6$~mol\'eculas~cm$^{-3}$ y que
          $\tau_{\ce{OH}} = \SI{10.6}{\hour}$.
          \textit{Exprese el resultado en cm$^3$~mol\'ecula$^{-1}$~s$^{-1}$.}

    \item A partir de sus resultados, explique por qu\'e los terpenos
          bios\'{\i}nteticos (isopreno, $\beta$-pineno) son especialmente
          importantes en la formaci\'on de aerosoles org\'anicos
          secundarios en atm\'osferas forestales.
\end{enumerate}

\textit{[Respuesta parcial:
(b) $\tau = \SI{10.6}{\hour} = 38160$~s;
    $k_{\ce{OH}} = 1/(38160 \times 10^6)
    \approx 2.6\times10^{-11}$~cm$^3$~mol\'ecula$^{-1}$~s$^{-1}$]}

% ----- Ejercicio 3: Química del SO2 ----------------------------------------
\item La oxidaci\'on del \ce{SO2} en la atm\'osfera sigue varias rutas.

\begin{enumerate}[label=\alph*)]
    \item Escriba la secuencia completa de cuatro reacciones que
          convierten \ce{SO2} en \ce{H2SO4} v\'{\i}a reacci\'on
          en fase gaseosa con el radical \ce{OH^.}, tal como se
          describe en el cap\'{\i}tulo. Identifique en qu\'e paso se
          regenera el radical \ce{HO2^.}.

    \item Se establece que el tiempo de vida del \ce{SO2} frente a
          la reacci\'on con \ce{OH^.} es de aproximadamente una semana.
          Si la concentraci\'on de \ce{OH^.} es
          $[\ce{OH}] = 1\times10^6$~mol\'eculas~cm$^{-3}$,
          estime la constante de velocidad de segundo orden
          $k_{\ce{OH},\ce{SO2}}$.
          \textit{Use $\tau \approx 7$~d.}

    \item El DMS (\ce{CH3SCH3}) es el mayor contribuidor natural
          al flujo global de azufre. Describa:
          (i)~su fuente biol\'ogica principal,
          (ii)~los dos oxidantes que controlan su vida en atm\'osferas
          marinas, y
          (iii)~por qu\'e su remoci\'on presenta un ciclo diurno.

    \item El \ce{H2S} tiene un tiempo de vida de $\sim$70~h frente
          a \ce{OH^.}. Escriba la reacci\'on inicial de remoci\'on
          y el producto final al que eventualmente conduce seg\'un
          el cap\'{\i}tulo.
\end{enumerate}

\textit{[Respuesta parcial:
(b) $\tau = 7 \times 86400 = 604800$~s;
    $k = 1/(604800 \times 10^6)
    \approx 1.6\times10^{-12}$~cm$^3$~mol\'ecula$^{-1}$~s$^{-1}$]}

% ----- Ejercicio 4: Estado fotoestacionario del ozono troposférico ----------
\item El sistema \ce{NO2}/\ce{NO}/\ce{O3} en la troposfera alcanza un
estado fotoestacionario descrito por las reacciones \ref{RN1}--\ref{RN3}
del cap\'{\i}tulo. La concentraci\'on de ozono en el estado estacionario est\'a
dada por:
\[
    [\ce{O3}] = \frac{k_1[\ce{NO2}]}{k_3[\ce{NO}]}
\]
y, cuando $[\ce{O3}]_0 = [\ce{NO}]_0 = 0$, por la expresi\'on
\ref{RN10}:
\[
    [\ce{O3}] = \frac{1}{2}\left[\left(\frac{k_1}{k_3}\right)^2
    + \frac{4k_1}{k_3}[\ce{NO2}]_0\right]^{1/2}
    - \frac{k_1}{2k_3}
\]

\begin{enumerate}[label=\alph*)]
    \item Usando $k_1/k_3 = 10$~ppb y $[\ce{NO2}]_0 = 100$~ppb,
          calcule $[\ce{O3}]$ en el estado estacionario y verifique
          que coincide con el valor del \textbf{Cuadro~\ref{ozono_foto}}.

    \item Repita el c\'alculo con $[\ce{NO2}]_0 = 1000$~ppb y compare
          con el valor tabulado.

    \item Aplique la aproximaci\'on de estado seudo-estable (PSSA) al
          \'atomo de ox\'{\i}geno \ce{O^.} para deducir la
          expresi\'on~\ref{RN7}:
          \[
              [\ce{O}] = \frac{k_1[\ce{NO2}]}{k_2[\ce{O2}][\ce{M}]}
          \]
          Explique f\'{\i}sicamente por qu\'e esta aproximaci\'on es
          v\'alida para una especie tan reactiva como \ce{O^.}.

    \item El cap\'{\i}tulo concluye que las concentraciones de \ce{O3}
          observadas en atm\'osferas urbanas (hasta 120~ppb) son
          mayores que las calculadas con s\'olo las reacciones
          \ref{RN1}--\ref{RN3}. ?`Qu\'e tipo de compuestos y
          reacciones adicionales son necesarios para explicar este
          exceso?
\end{enumerate}

\textit{[Respuesta parcial:
(a) $[\ce{O3}] = \tfrac{1}{2}[(10^2 + 4\times10\times100)^{1/2}] - 5
    = \tfrac{1}{2}[2300^{1/2}] - 5
    \approx \tfrac{1}{2}(47.96) - 5 \approx 19\ldots\approx 27$~ppb $\checkmark$;
(b) $[\ce{O3}] \approx 95$~ppb $\checkmark$]}

% ----- Ejercicio 5: Aminas y compuestos de carbono -------------------------
\item Los compuestos nitrogenados y carbonados desempe\~nan roles clave
en la qu\'{\i}mica atmosf\'erica.

\begin{enumerate}[label=\alph*)]
    \item Escriba las tres reacciones de aminas con el radical \ce{OH^.}
          presentadas en el cap\'{\i}tulo (para \ce{NH3},
          \ce{CH3NH2} y \ce{(CH3)2NH}) e identifique el tipo de
          producto radical generado en cada caso.

    \item El amon\'{\i}aco (\ce{NH3}) reacciona preferentemente con
          \'acidos atmosf\'ericos antes que con \ce{OH^.}.
          Escriba la reacci\'on con el \ce{HNO3} gaseoso e identifique
          el producto s\'olido formado. ?`Qu\'e relevancia tiene esta
          reacci\'on para la formaci\'on de PM$_{2.5}$?

    \item Describa las tres rutas de formaci\'on de radicales org\'anicos
          a partir de compuestos de carbono seg\'un el cap\'{\i}tulo
          (ox\'{\i}geno singelete, radical hidroxilo y ozono),
          escribiendo la reacci\'on general de cada una
          usando \ce{RH} como hidrocarburo gen\'erico.

    \item El PAN (nitrato de peroxiacetil) es un componente caracter\'{\i}stico
          del esmog fotoqu\'{\i}mico. Describa la secuencia de tres
          pasos que lleva desde un aldeh\'{\i}do hasta el PAN,
          pasando por el radical peroxiacilo y su reacci\'on final con
          \ce{NO2}.
\end{enumerate}

% ----- Ejercicio 6: Halógenos y destrucción catalítica ---------------------
\item Los compuestos hal\'ogenos participan en ciclos catal\'{\i}ticos
de destrucci\'on de ozono.

\begin{enumerate}[label=\alph*)]
    \item La reacci\'on del \'atomo de Cl con hidrocarburos (reacci\'on
          \ref{RH1}) compite con su oxidaci\'on por \ce{O3}
          (reacci\'on \ref{RH2}). Seg\'un el cap\'{\i}tulo, la fracci\'on
          que reacciona con \ce{O3} en lugar de \ce{RH} es:
          F $\approx$ 0\% (F), 50\% (Cl), 99\% (Br) y 100\% (I).
          Explique cualitativamente esta tendencia en t\'erminos de la
          reactividad decreciente de los hal\'ogenos hacia compuestos
          que contienen H (HX).

    \item La reacci\'on del \ce{CH3Cl} con \ce{OH^.} libera finalmente
          un \'atomo de \ce{Cl^.}. Escriba la secuencia completa de
          tres reacciones indicadas en el cap\'{\i}tulo:
          \ce{CH3Cl + OH^. ->}, \ce{^.CH2Cl + O2 ->} y
          \ce{CH2ClOO^. + NO ->}.

    \item Describa c\'omo el \ce{HCl} generado en la estratosfera
          puede ser devuelto al ciclo activo del cloro (\ce{ClO_$x$})
          mediante la reacci\'on con el radical \ce{OH^.} (reacci\'on
          \ref{RH3}).

    \item Explique por qu\'e las constantes de velocidad de los \'atomos
          de Cl con hidrocarburos son significativamente mayores que las
          del \ce{OH^.} con los mismos compuestos, y qu\'e consecuencia
          tiene esto para la remoci\'on de hidrocarburos en zonas
          costeras con aerosol salino marino.
\end{enumerate}

% ----- Ejercicio 7: Balance de ozono estratosférico (Chapman) --------------
\item El balance natural de ozono en la estratosfera se describe mediante
las reacciones de Chapman.

\begin{enumerate}[label=\alph*)]
    \item Escriba las cuatro reacciones del mecanismo de Chapman:
          dos de formaci\'on (reacciones \ref{OE1} y \ref{OE2})
          y dos de destrucci\'on (reacci\'on \ref{OE3} y su reacci\'on
          inversa de regeneraci\'on). Indique las longitudes de onda
          de los fotones involucrados en cada fotolisis.

    \item Usando la analog\'{\i}a del ``balde que gotea'' descrita en
          el cap\'{\i}tulo, explique qu\'e condici\'on debe cumplirse para
          que el nivel de ozono permanezca constante, y qu\'e ocurre
          cuando se introducen compuestos adicionales como los CFC.

    \item El cap\'{\i}tulo indica que la concentraci\'on promedio de ozono
          es 300~DU, equivalente a una capa de \SI{3}{\milli\metre}.
          Si en el ``agujero'' ant\'artico cae a 100~DU:
          \begin{enumerate}[label=\roman*)]
              \item ?`Cu\'antos mil\'{\i}metros de espesor
                    representan 100~DU?
              \item Usando la definici\'on de que 1~DU $\equiv$
                    $2.69\times10^{16}$~mol\'eculas~cm$^{-2}$,
                    calcule el n\'umero total de mol\'eculas de
                    \ce{O3} en una columna de 1~cm$^2$ de base
                    para 300~DU y para 100~DU.
          \end{enumerate}

    \item La columna total de aire comprimido a STP tiene un espesor
          de \SI{8}{\kilo\metre}. Calcule qu\'e porcentaje del espesor
          total de la atm\'osfera representa la capa de ozono de 300~DU
          (\SI{3}{\milli\metre}).
\end{enumerate}

\textit{[Respuesta parcial:
(c-i) $100\,\text{DU} \rightarrow \SI{1}{\milli\metre}$;
(c-ii) $300\,\text{DU}: N = 300 \times 2.69\times10^{16}
      = 8.07\times10^{18}$~mol\'eculas~cm$^{-2}$;
      $100\,\text{DU}: N = 2.69\times10^{18}$~mol\'eculas~cm$^{-2}$;
(d) $3\,\text{mm}/8\,000\,000\,\text{mm} = 3.75\times10^{-7}$, es
    decir, $\sim 0.000038\%$ del espesor total]}

% ----- Ejercicio 8: Ciclos catalíticos de pérdida de ozono -----------------
\item Los ciclos catal\'{\i}ticos \ce{HO_$x$}, \ce{NO_$x$} y
\ce{ClO_$x$} son responsables de la p\'erdida acelerada de ozono
en la estratosfera.

\begin{enumerate}[label=\alph*)]
    \item Para el ciclo \ce{HO_$x$}, escriba las dos reacciones
          propagadoras (\ce{OH + O3} y \ce{HO2 + O3}) y la reacci\'on
          global neta. ?`Cu\'al es el cat\'alizador del ciclo?
          ?`C\'omo se termina el ciclo?

    \item Para el ciclo catal\'{\i}tico del \ce{NO_$x$}
          (reacciones \ref{R10}--\ref{R11}), escriba las dos
          reacciones y la reacci\'on neta. Calcule cu\'antas
          mol\'eculas de \ce{O_$x$} (equivalente a \ce{O3})
          se consumen por cada ciclo completo.

    \item El ciclo \ce{ClO_$x$} se inicia con la fotolisis del CFC-12
          (\ce{CF2Cl2}). Escriba:
          (i)~la reacci\'on de fotolisis del CFC-12,
          (ii)~las dos reacciones propagadoras del ciclo
          \ce{Cl}/\ce{ClO}, y
          (iii)~la reacci\'on neta global.

    \item Compare los tres ciclos catal\'{\i}ticos (\ce{HO_$x$},
          \ce{NO_$x$}, \ce{ClO_$x$}) llenando la siguiente tabla:

{%
\centering
\begin{tabular}{lccc}
\toprule
\textbf{Caracter\'{\i}stica} & \ce{HO_$x$} & \ce{NO_$x$} & \ce{ClO_$x$} \\
\midrule
Radical cat. & & & \\
Fuente principal & & & \\
Reacci\'on de terminaci\'on & & & \\
Reservorio inactivo & & & \\
\bottomrule
\end{tabular}
\par
}%
\end{enumerate}

% ----- Ejercicio 9: PSCs y agujero de ozono --------------------------------
\item Las nubes estratosf\'ericas polares (PSC) desempe\~nan un papel
fundamental en el agujero de ozono ant\'artico.

\begin{enumerate}[label=\alph*)]
    \item El cap\'{\i}tulo indica que las PSC se forman alrededor de
          \SI{-72}{\celsius}. Convierta esta temperatura a Kelvin
          y explique por qu\'e las PSC son m\'as frecuentes sobre
          la Ant\'artida que sobre el \'Artico.

    \item En la superficie de las PSC los gases de ``dep\'osito''
          \ce{HCl} y \ce{ClONO2} se activan. Escriba la reacci\'on
          heterog\'enea de \ce{N2O5} con \ce{NaCl} en aerosol salino
          que produce una especie activa de cloro (\ce{XNO2}) y una sal
          s\'olida, seg\'un el cap\'{\i}tulo.

    \item La reacci\'on \ce{N2O + O(^1D) -> 2NO} en la estratosfera
          es responsable del 5\% de la p\'erdida del \ce{N2O}.
          Si la tasa total de p\'erdida de \ce{N2O} es
          $F = \SI{1.0e10}{\kilo\gram}/$a\~no, ?`cu\'anta masa
          (en kg/a\~no) se convierte a NO por esta v\'{\i}a?
          ?`Qu\'e ocurre con el 95\% restante?

    \item Molina y Rowland recibieron el Premio Nobel de Qu\'{\i}mica
          en 1995 junto con Paul Crutzen. Explique la contribuci\'on
          espec\'{\i}fica de cada uno en la comprensi\'on de la
          destrucci\'on del ozono estratosf\'erico, con base en la
          informaci\'on del cap\'{\i}tulo.
\end{enumerate}

\textit{[Respuesta parcial:
(a) \SI{-72}{\celsius} $= \SI{201}{\kelvin}$;
(c) $5\%$ de $10^{10}$ kg/a\~no $= 5\times10^8$~kg/a\~no;
    el 95\% se convierte en \ce{N2} (inerte) por fotolisis y  oxidaci\'on con \ce{O(^1D)}]}

% ----- Ejercicio 10 (avanzado): Ciclo 2 de ClO y ecuación de continuidad ---
\item \textbf{[Avanzado]} El cap\'{\i}tulo describe tres ciclos de
destrucci\'on de ozono con hal\'ogenos (\textbf{Cuadro~\ref{desozo}}).

\begin{enumerate}[label=\alph*)]
    \item Para el \textbf{Ciclo 2} (\ce{ClO + ClO -> (ClO)2 -> \ldots}),
          escriba los cuatro pasos elementales y la reacci\'on global
          neta. Identifique cu\'al paso requiere un fot\'on ($h\nu$)
          y cu\'al produce el \'atomo de Cl libre necesario para
          continuar el ciclo.

    \item Para el \textbf{Ciclo 3} (que involucra tanto Cl como Br),
          escriba los cuatro pasos y la reacci\'on neta.
          Compare la eficiencia del Ciclo 3 con la del Ciclo 1
          en t\'erminos de la raz\'on de mol\'eculas de \ce{O3}
          destruidas por \'atomo de hal\'ogeno.

    \item La ecuaci\'on de continuidad del ozono estratosf\'erico es:
          \[
              \frac{\partial[\ce{O3}]}{\partial t}
              = P - L - \nabla\cdot(\vec{v}[\ce{O3}])
          \]
          Identifique cada t\'ermino f\'{\i}sicamente e indique
          qu\'e proceso representa cada uno:
          \begin{enumerate}[label=\roman*)]
              \item el t\'ermino $P$,
              \item el t\'ermino $L$,
              \item el t\'ermino $\nabla\cdot(\vec{v}[\ce{O3}])$.
          \end{enumerate}
          ?`En qu\'e regi\'on de la Tierra domina el transporte
          advectivo sobre la producci\'on qu\'{\i}mica?

    \item La circulaci\'on de Brewer-Dobson transporta \ce{O3} desde
          los tr\'opicos hacia los polos. Explique por qu\'e, a pesar
          de que el \ce{O3} se produce principalmente en la estratosfera
          tropical, las m\'aximas concentraciones se observan en
          latitudes altas (salvo en la Ant\'artida en primavera austral).
\end{enumerate}

\end{ejercicios}

\end{bloque_ejercicios}
% =============================================================================
% FIN DE ejercicios_cap08.tex
% =============================================================================


